% p251_1; source PDF page 251; manually verified editable table.
\begin{tikzpicture}[x=1cm,y=-1cm,inner sep=0pt,outer sep=0pt]
\path[use as bounding box] (0,0) rectangle (16.000,5.100);
\draw[line width=.4pt] (0.000,0.000) rectangle (2.200,0.850);
\node[align=center,text width=1.940cm,font={\linespread{1}\fontsize{10}{13.50}\selectfont\bfseries}] at (1.100,0.425) {对比};
\draw[line width=.4pt] (2.200,0.000) rectangle (6.000,0.850);
\node[align=center,text width=3.540cm,font={\linespread{1}\fontsize{10}{13.50}\selectfont\bfseries}] at (4.100,0.425) {虐待罪};
\draw[line width=.4pt] (6.000,0.000) rectangle (16.000,0.850);
\node[align=center,text width=9.740cm,font={\linespread{1}\fontsize{10}{13.50}\selectfont\bfseries}] at (11.000,0.425) {虐待被监护、看护人罪};
\draw[line width=.4pt] (0.000,0.850) rectangle (2.200,2.550);
\node[align=left,text width=1.940cm,font={\linespread{1}\fontsize{10}{13.50}\selectfont}] at (1.100,1.700) {行为对象};
\draw[line width=.4pt] (2.200,0.850) rectangle (6.000,2.550);
\node[align=left,text width=3.540cm,font={\linespread{1}\fontsize{10}{13.50}\selectfont}] at (4.100,1.700) {家庭成员};
\draw[line width=.4pt] (6.000,0.850) rectangle (16.000,2.550);
\node[align=left,text width=9.740cm,font={\linespread{1}\fontsize{10}{13.50}\selectfont}] at (11.000,1.700) {被监护、看护人（两罪的行为对象是A与B的关系，而非A与A+B。因此，两罪不是法条竞合关系，而是想象竞合关系）};
\draw[line width=.4pt] (0.000,2.550) rectangle (2.200,3.400);
\node[align=left,text width=1.940cm,font={\linespread{1}\fontsize{10}{13.50}\selectfont}] at (1.100,2.975) {行为主体};
\draw[line width=.4pt] (2.200,2.550) rectangle (6.000,3.400);
\node[align=left,text width=3.540cm,font={\linespread{1}\fontsize{10}{13.50}\selectfont}] at (4.100,2.975) {只能由自然人构成};
\draw[line width=.4pt] (6.000,2.550) rectangle (16.000,3.400);
\node[align=left,text width=9.740cm,font={\linespread{1}\fontsize{10}{13.50}\selectfont}] at (11.000,2.975) {自然人和单位};
\draw[line width=.4pt] (0.000,3.400) rectangle (2.200,4.250);
\node[align=left,text width=1.940cm,font={\linespread{1}\fontsize{10}{13.50}\selectfont}] at (1.100,3.825) {案件性质};
\draw[line width=.4pt] (2.200,3.400) rectangle (6.000,4.250);
\node[align=left,text width=3.540cm,font={\linespread{1}\fontsize{10}{13.50}\selectfont}] at (4.100,3.825) {告诉才处理，有例外};
\draw[line width=.4pt] (6.000,3.400) rectangle (16.000,4.250);
\node[align=left,text width=9.740cm,font={\linespread{1}\fontsize{10}{13.50}\selectfont}] at (11.000,3.825) {公诉案件};
\draw[line width=.4pt] (0.000,4.250) rectangle (2.200,5.100);
\node[align=left,text width=1.940cm,font={\linespread{1}\fontsize{10}{13.50}\selectfont}] at (1.100,4.675) {结果加重犯};
\draw[line width=.4pt] (2.200,4.250) rectangle (6.000,5.100);
\node[align=left,text width=3.540cm,font={\linespread{1}\fontsize{10}{13.50}\selectfont}] at (4.100,4.675) {有结果加重犯};
\draw[line width=.4pt] (6.000,4.250) rectangle (16.000,5.100);
\node[align=left,text width=9.740cm,font={\linespread{1}\fontsize{10}{13.50}\selectfont}] at (11.000,4.675) {无结果加重犯};
\end{tikzpicture}
